% Replacement text for the supplied manuscript; not a standalone document.
% Replace Table 1 entries: T=5000; i_incentive=2.0.
% Replace Table 2 entries: T=5000; input duration=10.5 h; beta=13;
% timeout=3920; outflow target=65000; F=1000000 RL control decisions.

\add{The PeMS inputs span 03:00--13:30 local time, with boundary-flow scales of 0.6, 0.8, and 1.0 for three, four, and five lanes, respectively. Native training uses a budget of one million RL control decisions per seed. The physical evaluation includes both entrance queues and continues until complete clearance, subject to a safety cap of 28,800 CTM steps.}

\add{Figure~12(c) compares ten weekday dates, while Figure~12(d) reports seven day-of-week groups using fourteen available dates. Each date includes thirty paired demand scenarios. The intervals are conditional on the available dates and trained models; these comparisons alone do not establish unseen-date generalization.}

\add{Figure~5 presents an illustrative stress test with temporally correlated CAV arrivals. The mainline and on-ramp have separate randomly generated 60-minute composition regimes, within which CAV fractions fluctuate at every 15-second CTM step. The vehicle-weighted CAV fraction across the complete offered demand is 0.4. Both controllers receive identical demand and initial conditions; the trained PPO, five-term reward, and traffic model are unchanged. In this example, ADDLC achieves 5,368 more downstream departures by the end of the 10.5-hour demand window and reduces complete TTS, including both entrance queues until clearance, from 556,285 to 472,896 vehicle-hours (14.99\%). This case illustrates the response to correlated composition changes and is separate from the independent-step composition distribution used in the main training and reward-ablation experiments; it does not establish an advantage for every demand realization.}

\add{The lane-width comparison uses the same saved cohort of one hundred cases for each road width. Both boundary flows are scaled by 0.8 for four lanes and 1.0 for five lanes. The four-lane results use the subsequently retrained policies, while the five-lane results remain unchanged. Relative to direct transfer, lane-specific retraining reduces complete TTS by 0.87\% and 15.88\%, and mean entry waiting time by 2.75\% and 13.64\%, for four and five lanes, respectively. Fixed-window downstream throughput increases by 0.25\% and 4.29\%. The four-lane episode rewards are practically similar (4.728 for retraining and 4.730 for transfer). These results describe the saved demonstration cohorts; they do not establish population-wide superiority or isolate the effect of road width.}
